\chapter{Análisis semántico y sistema de tipos}
\label{ch:analisis-semantico}

El análisis semántico de Qalc tiene como objetivo dotar a cada nodo del AST de un tipo completamente resuelto, verificando que el programa es coherente desde el punto de vista tipológico antes de que ninguna fase posterior lo transforme en código ejecutable. Este proceso se apoya en dos mecanismos que se ejecutan de forma coordinada, el ámbito léxico, que resuelve a qué declaración se refiere cada identificador, y el inferidor de tipos basado en el algoritmo de Hindley-Milner, que asigna un tipo a cada expresión a partir de esas resoluciones. Este capítulo describe primero la resolución del ámbito léxico, después el algoritmo de inferencia con sus tres fases, la representación interna de los tipos durante la inferencia, el mecanismo de unificación en detalle y la sobrecarga de funciones discriminada por tipo.

\section{Ámbito léxico y declaraciones de nivel superior}
\label{sec:ambito-lexico}

El primer paso del análisis semántico consiste en resolver el ámbito léxico de las declaraciones, de modo que la fase de inferencia pueda tratar cada identificador exactamente como el símbolo al que se refiere. El módulo \lstinline|qalc-ast/scope.rs| implementa el ámbito léxico del compilador. Un \lstinline|Scope| es un mapa de nombres a vectores de \lstinline|SymbolInfo|, con un puntero opcional al ámbito padre para implementar el encadenamiento léxico:

\begin{lstlisting}[language={}, caption=Definición del tipo Scope para el ámbito léxico]
pub struct Scope {
    parent: Option<Rc<Scope>>,
    symbols: RefCell<HashMap<String, Vec<SymbolInfo>>>,
}
\end{lstlisting}

Hay dos operaciones de inserción con semánticas distintas. La primera, \lstinline|define(&self, symbol)|, registra un nombre nuevo y, si el nombre ya existe en el ámbito local, devuelve \lstinline|false| sin sobreescribirlo (se usa para parámetros y variables \lstinline|let|, donde la redefinición en el mismo ámbito es un error). La segunda, \lstinline|define_decl(&self, symbol)|, añade la definición al vector del nombre sin rechazar las entradas previas, lo que permite acumular varias definiciones para el mismo nombre (se usa para declaraciones de función de nivel superior, donde el mismo nombre puede tener varias definiciones con firmas de tipo distintas; la resolución de la sobrecarga se describe en la \autoref{sec:overloading}).

Al inicio de la inferencia, el \lstinline|HMInferrer| realiza dos operaciones preparatorias sobre el ámbito global antes de visitar ninguna declaración del programa.

En primer lugar, inyecta todas las funciones predefinidas del crate \lstinline|qalc-builtins|. Cada \lstinline|BuiltinDef| contiene el nombre de la función y su tipo completo expresado como \lstinline|TypeWithUnknow|, de modo que el inferidor puede resolver referencias a built-ins exactamente igual que referencias a funciones definidas por el usuario.

En segundo lugar, recorre todas las declaraciones del programa y registra sus firmas de tipo en el ámbito global antes de inferir sus cuerpos. Este paso de predeclaración es fundamental para que las funciones puedan referenciarse mutuamente sin que el orden de aparición en el fichero fuente afecte a la corrección de la inferencia.

\section{El algoritmo de Hindley-Milner}
\label{sec:hm}

Una vez resuelto el ámbito léxico de las declaraciones, Qalc infiere el tipo de cada expresión mediante el algoritmo de Hindley-Milner, que resuelve el problema de la inferencia de tipos en un lenguaje funcional sin necesidad de que el programador anote los tipos de cada expresión. A partir de las anotaciones de las declaraciones de nivel superior y de las reglas del sistema de tipos, el compilador deduce el tipo de toda subexpresión.

El nombre del algoritmo combina las aportaciones de dos autores. Hindley estudió en 1969 la asignación de tipos principales en el marco de la lógica combinatoria~\cite{hindley_1969}, y Milner reformuló esa idea en 1978 como un algoritmo de inferencia práctico, conocido como algoritmo W~\cite{milner_1978}, pensado para el lenguaje ML. Damas y Milner~\cite{damas_milner_1982} extendieron después el algoritmo para admitir polimorfismo en las declaraciones \lstinline|let|, dando lugar a la versión que se usa habitualmente hoy en día y que sirve de referencia para la implementación de Qalc.

El algoritmo de Hindley-Milner es la base de la inferencia de tipos en toda la familia de lenguajes ML (del inglés \textit{Meta Language}), que incluye Standard ML, OCaml y F#, y en Haskell, que lo extiende con clases de tipos para resolver la sobrecarga. Lenguajes más recientes como Elm o PureScript también lo adoptan sin modificaciones sustanciales. Otros lenguajes con un sistema de tipos más flexible, como Rust~\cite{matsakis_rust_2014}, Swift, Kotlin o TypeScript, usan variantes locales o bidireccionales inspiradas en Hindley-Milner que infieren el tipo de expresiones concretas sin llegar al polimorfismo global de la versión completa del algoritmo. Qalc adopta Hindley-Milner porque permite prescindir de anotaciones de tipo en la mayoría de las expresiones sin renunciar a la seguridad de tipos, una propiedad especialmente valiosa en un lenguaje pensado para describir circuitos cuánticos de forma concisa.

El algoritmo se estructura en tres fases que se ejecutan secuencialmente~\cite{pierce_tapl_2002}:

\begin{enumerate}
  \item Generación de restricciones. Se recorre el AST y a cada subexpresión se le asigna una variable de tipo fresca. A medida que se visitan los nodos, se generan restricciones de tipo, es decir, aserciones de la forma $\tau_1 = \tau_2$ que afirman que dos tipos deben ser iguales. Por ejemplo, al analizar la aplicación $f\;a$, se genera la restricción $\tau_f = \tau_a \to \tau_{resultado}$.
  \item Unificación. Las restricciones generadas se resuelven mediante el algoritmo de unificación de Robinson~\cite{robinson_1965}. El resultado es una sustitución, es decir, una función finita que asigna a cada variable de tipo fresca el tipo concreto (o el tipo estructurado) que la resuelve.
  \item Aplicación de la sustitución. La sustitución se aplica a todo el AST, reemplazando cada variable de tipo por su valor en la sustitución. El resultado es el \lstinline|TypedAST|, en el que cada nodo lleva su tipo completamente resuelto.
\end{enumerate}

En Qalc, estas tres fases se implementan en los módulos \lstinline|infer.rs|, \lstinline|unify.rs| y \lstinline|apply.rs| del crate \lstinline|qalc-analysis/hm_inferrer|, respectivamente.

Para ilustrar cómo cooperan las tres fases sobre un caso concreto, considérese la siguiente declaración.

\begin{lstlisting}[language={}, caption=Ejemplo de declaración a inferir mediante Hindley-Milner]
incrementa :: i32 -> i32
incrementa = \x -> x + 1
\end{lstlisting}

En la fase de generación de restricciones (\autoref{sec:constraint-gen}), la lambda asigna una variable fresca $\tau_x$ al parámetro \lstinline|x| y visita el cuerpo \lstinline|x + 1|. Como \lstinline|+| es una función predefinida de tipo \lstinline|i32 -> i32 -> i32|, la aplicación genera las restricciones $\tau_x = \mathtt{i32}$ y $\tau_{cuerpo} = \mathtt{i32}$, y la anotación de la declaración añade además la restricción $\tau_x \to \tau_{cuerpo} = \mathtt{i32} \to \mathtt{i32}$.

En la fase de unificación (\autoref{sec:unificacion}), estas restricciones se resuelven en una única sustitución que asigna $\tau_x \mapsto \mathtt{i32}$ y $\tau_{cuerpo} \mapsto \mathtt{i32}$, sin que se detecte ningún conflicto entre las variables y los tipos concretos involucrados.

En la fase de aplicación de la sustitución, esa sustitución se propaga a cada nodo de la lambda, de modo que el \lstinline|TypedAST| resultante registra que \lstinline|x| tiene tipo \lstinline|i32| y que \lstinline|incrementa| tiene tipo \lstinline|i32 -> i32|, coherente con la anotación original de la declaración.

\section{Representación de tipos durante la inferencia}

El crate \lstinline|qalc-ast| define dos enumeraciones para representar los tipos en distintos momentos del pipeline.

\begin{itemize}
  \item \texttt{TypeWithUnknow.} Es la representación usada durante la inferencia. Sus variantes son las siguientes: \lstinline|Unknow(usize)| (una variable de tipo fresca, identificada por un entero único generado con \lstinline|gen_id()|, que representa tipos aún no resueltos), \lstinline|Builtin \{ id, params, linear \}| (un tipo conocido, como \lstinline|i32|, \lstinline|bool| o \lstinline|qbool|, donde el campo \lstinline|linear| distingue los tipos cuánticos de los clásicos), \lstinline|Arrow(a, b)| (el tipo función $a \to b$), \lstinline|Tuple(ts)| (el tipo tupla $(t_1, t_2, \ldots, t_n)$) y \lstinline|Slice(t)| (el tipo slice $[t]$).

  \item \texttt{Type.} Es la representación final, producida una vez que la unificación ha resuelto todas las variables. Es isomorfa a \lstinline|TypeWithUnknow| salvo por la ausencia de la variante \lstinline|Unknow|: si la unificación no pudo resolver alguna variable, la conversión \lstinline|TryFrom<TypeWithUnknow> for Type| falla con un error de compilación que pide al usuario una anotación de tipo explícita.
\end{itemize}

La separación entre \lstinline|TypeWithUnknow| y \lstinline|Type| garantiza, a nivel del sistema de tipos de Rust, que ningún nodo del \lstinline|TypedAST| final pueda contener una variable sin resolver.

\section{Generación de restricciones de tipo}
\label{sec:constraint-gen}

La función \lstinline|infer_expr| recorre el AST normalizado en profundidad y produce un par \lstinline|(Expression<Inferred>, Vec<TypeConstraint>)|: el nodo anotado con su variable de tipo y el conjunto de restricciones que impone la estructura de ese nodo.

Las reglas más relevantes son las siguientes:

\begin{itemize}
  \item \textit{Literales.} Los literales enteros y flotantes producen directamente el tipo concreto determinado por su sufijo: \lstinline|42i32| produce \lstinline|Builtin("i32")|, \lstinline|3.14f64| produce \lstinline|Builtin("f64")|. Los literales sin sufijo explícito se tipan por defecto como \lstinline|i32| y \lstinline|f64| respectivamente.

  \item \textit{Identificadores.} La resolución de un identificador consulta el ámbito léxico en vigor. Si el nombre se encuentra, el nodo hereda el tipo registrado en el símbolo. Si no existe, se emite un error de compilación con el span del identificador.

  \item \textit{Aplicación de función} (\texttt{Transform \{ callee, param \}}). Se genera una variable fresca $\tau_{ret}$ para el tipo de retorno y se añade la restricción $\tau_{callee} = \tau_{param} \to \tau_{ret}$. El tipo del nodo de aplicación es $\tau_{ret}$.

  \item \textit{Lambda.} Se crea un ámbito nuevo, se registra cada parámetro con una variable fresca y se infiere el cuerpo. El tipo de la lambda se construye encadenando los tipos de los parámetros con el tipo del cuerpo, tal y como muestra la expresión
  \[\tau_{lambda} = \tau_{p_1} \to \tau_{p_2} \to \cdots \to \tau_{cuerpo}\]

  \item \textit{Bloque} \texttt{let}. Cada enlace $p = e$ genera las restricciones de $e$ y añade la restricción $\tau_e = \tau_p$, ligando el tipo de la expresión al tipo del patrón. Después se infiere la expresión de valor en el ámbito extendido. El tipo del bloque es el tipo de la expresión de valor.

  \item \textit{If-else.} Se genera una variable fresca $\tau_{ret}$ y se añaden las restricciones $\tau_{ret} = \tau_{then}$ y $\tau_{ret} = \tau_{else}$, exigiendo que ambas ramas tengan el mismo tipo. El tipo del test no se restringe a ningún tipo concreto en esta fase.

  \item \textit{Tuplas.} El tipo de una tupla $(e_1, \ldots, e_n)$ es \lstinline|Tuple|$(\tau_{e_1}, \ldots, \tau_{e_n})$, donde cada $\tau_{e_i}$ es el tipo inferido de la subexpresión correspondiente.

  \item \textit{Slices.} Se genera una variable fresca $\tau_{elem}$ y se añade la restricción $\tau_{elem} = \tau_{e_i}$ para cada elemento, exigiendo homogeneidad del slice.
\end{itemize}

La función \lstinline|infer_decl| maneja las declaraciones de función. Cada declaración tiene una anotación de tipo explícita escrita por el programador. El inferidor convierte esa anotación en un \lstinline|TypeWithUnknow| mediante \lstinline|type_expr_to_infer_type| y añade la restricción:
\[\tau_{inferida} = \tau_{anotada}\]

Esto obliga al cuerpo de la función a ser coherente con la firma declarada. Si la inferencia produce un tipo incompatible, el error se reporta apuntando al span de la declaración.

\section{Unificación}
\label{sec:unificacion}

La unificación resuelve el conjunto de restricciones generado en la fase anterior. El algoritmo implementado en \lstinline|unify.rs| es el de Robinson~\cite{robinson_1965}, en su variante con verificación de ocurrencias (\textit{occurs check}).

El algoritmo mantiene una sustitución parcial $S$ (un mapa de identificadores de variables a tipos). Para cada restricción $\tau_1 = \tau_2$, primero aplica $S$ a ambos lados (reemplazando las variables ya resueltas) y después llama a \lstinline|unify_one|:

\begin{itemize}
  \item Si $\tau_1 = \tau_2$, no hay nada que hacer.
  \item Si alguno de los dos lados es una variable fresca $?n$, se extiende $S$ con $n \mapsto \tau_{otro}$, siempre que $n$ no aparezca en $\tau_{otro}$ (occurs check). La verificación de ocurrencias previene la generación de tipos circulares infinitos como $\tau = \tau \to \tau$.
  \item Si ambos son tipos estructurados del mismo constructor (\lstinline|Arrow|, \lstinline|Tuple|, \lstinline|Slice|, \lstinline|Builtin|), se unifican recursivamente sus subcomponentes.
  \item En cualquier otro caso, los tipos son incompatibles y se emite un error de compilación con el span de la restricción original, indicando qué tipos no pudieron unificarse.
\end{itemize}

La implementación aplica la sustitución parcial acumulada antes de cada unificación, lo que permite que el resultado de resolver una restricción temprana informe de inmediato a las restricciones posteriores, sin necesidad de múltiples pasadas.

\section{Aplicación de la sustitución}

Una vez que \lstinline|unify| devuelve la sustitución final $S$, el módulo \lstinline|apply.rs| recorre el AST \lstinline|Inferred| de forma estructural, reemplazando cada \lstinline|TypeWithUnknow::Unknow(id)| por \lstinline|S(id)| en cada nodo. El resultado es el \lstinline|TypedAST|, en el que todas las variables de tipo han sido resueltas a tipos concretos y el AST puede convertirse al tipo final \lstinline|Typed|.

Si alguna variable permanece irresuelta tras la sustitución (por ejemplo, en una función sin anotación de tipo y con un polimorfismo no instanciado), la conversión \lstinline|TryFrom<TypeWithUnknow> for Type| falla y se reporta un error pidiendo al usuario que añada una anotación explícita.

\section{Sobrecarga de funciones discriminada por tipo}
\label{sec:overloading}

Una de las extensiones del sistema de tipos es la sobrecarga de funciones basada en el tipo del primer parámetro. En Qalc, una función puede tener varias definiciones con el mismo nombre siempre que difieran en el tipo del primer parámetro; el compilador elige la definición adecuada en tiempo de compilación consultando el tipo del primer argumento en el punto de llamada.

Este mecanismo permite expresar operaciones polimórficas sin necesidad de clases de tipos ni de polimorfismo paramétrico completo, manteniendo la inferencia de tipos decidible y eficiente. Por ejemplo, la función \lstinline|at|, que aplica una puerta a un cúbit concreto, puede estar definida dos veces con el mismo nombre:

\begin{lstlisting}[language={}, caption=Ejemplo de sobrecarga de función discriminada por tipo del primer parámetro]
at :: i32 -> (qcirc -> qcirc) -> qcirc -> qcirc
at :: [i32] -> (qcirc -> qcirc) -> qcirc -> qcirc
\end{lstlisting}

La sobrecarga está soportada a nivel de ámbito: \lstinline|define_decl| acumula todas las definiciones de un mismo nombre en un vector, y la operación \lstinline|resolve| recorre ese vector dado el tipo del primer argumento en el punto de llamada, y elige la definición cuyo primer tipo de parámetro es unificable con el tipo del argumento. Si ninguna definición es compatible, se emite un error de compilación indicando las firmas disponibles y el tipo que se intentó usar.

La discriminación ocurre tras la generación de restricciones para el primer argumento, cuando su tipo ya está disponible (posiblemente aún como variable fresca parcialmente restringida). En los casos en que el tipo del primer argumento no está suficientemente determinado en ese punto, el compilador difiere la elección hasta después de la unificación y resuelve la sobrecarga en una segunda pasada sobre el AST, manteniendo en el nodo de identificador una referencia al conjunto de candidatos hasta que la sustitución lo resuelve.

Esta estrategia de sobrecarga es deliberadamente más sencilla que los sistemas de clases de tipos de Haskell o los rasgos de Rust~\cite{matsakis_rust_2014}, ya que restringir la discriminación al tipo del primer parámetro garantiza que la resolución es siempre unívoca y que no es necesario un mecanismo de inferencia global adicional. El coste es una menor expresividad, que en el contexto de Qalc es aceptable dado que el objetivo del lenguaje es describir circuitos cuánticos, no implementar abstracciones de propósito general.

\section{Casos de uso}
\label{sec:casos-uso}

Los siguientes ejemplos recorren el proceso completo de análisis semántico descrito en este capítulo sobre fragmentos de código más cercanos al uso real del lenguaje que los ejemplos aislados de las secciones anteriores, mostrando cómo cooperan la resolución de ámbito, la inferencia de tipos, la unificación y la sobrecarga sobre un mismo programa.

\subsection{Inferencia sobre una operación cuántica compuesta}

Considérese la siguiente declaración, que aplica la puerta Hadamard a dos cúbits de entrada y devuelve ambos resultados en una tupla.

\begin{lstlisting}[language={}, caption=Declaración que combina un bloque let y una tupla sobre valores cuánticos]
prepara :: qbool -> qbool -> (qbool, qbool)
prepara a b = let x = H a in (x, H b)
\end{lstlisting}

La resolución de ámbito (\autoref{sec:ambito-lexico}) registra los parámetros \lstinline|a| y \lstinline|b| con sendas variables frescas $\tau_a$ y $\tau_b$, y resuelve el identificador \lstinline|H| contra la firma \lstinline|qbool -> qbool| que \lstinline|qalc-builtins| inyectó en el ámbito global.

Durante la generación de restricciones (\autoref{sec:constraint-gen}) se combinan cuatro de las reglas descritas en ese apartado. La aplicación \lstinline|H a| impone, según la regla de aplicación de función, la restricción $\mathtt{qbool} \to \mathtt{qbool} = \tau_a \to \tau_{ret_1}$, de la que se derivan $\tau_a = \mathtt{qbool}$ y $\tau_{ret_1} = \mathtt{qbool}$. La regla del bloque \lstinline|let| liga entonces $\tau_x = \tau_{ret_1}$. La segunda aplicación \lstinline|H b| genera de forma análoga $\tau_b = \mathtt{qbool}$ y una variable fresca $\tau_{ret_2} = \mathtt{qbool}$, y la regla de tuplas construye el tipo \lstinline|Tuple|$(\tau_x, \tau_{ret_2})$ para el resultado. Por último, la anotación de la declaración añade la restricción de que el tipo completo coincida con \lstinline|qbool -> qbool -> (qbool, qbool)|.

La unificación (\autoref{sec:unificacion}) resuelve todas las variables anteriores a \lstinline|qbool| sin ningún conflicto, y la sustitución confirma que \lstinline|prepara| tiene exactamente el tipo anotado. El carácter lineal de \lstinline|qbool| viaja con el tipo sin intervención adicional del inferidor, puesto que forma parte del propio campo \lstinline|linear| de la variante \lstinline|Builtin|; comprobar que \lstinline|a| y \lstinline|b| se consumen exactamente una vez es responsabilidad de la fase de comprobación de consumo lineal, posterior al aplanado y descrita en el \autoref{sec:linear-checker}.

\subsection{Resolución de sobrecarga en una llamada concreta}

Retomando la función \lstinline|at| introducida en el \autoref{sec:overloading}, considérense dos llamadas sobre el mismo circuito base.

\begin{lstlisting}[language={}, caption=Dos llamadas a la función sobrecargada at resueltas con distinta firma]
circuito1 = at 0 H base
circuito2 = at [0, 1, 2] H base
\end{lstlisting}

En la primera llamada, la generación de restricciones infiere que el literal \lstinline|0| tiene tipo \lstinline|i32|. Al consultar el vector de definiciones de \lstinline|at| en el ámbito, \lstinline|resolve| encuentra que \lstinline|i32| unifica con el primer parámetro de la firma \lstinline|i32 -> (qcirc -> qcirc) -> qcirc -> qcirc| y descarta la firma cuyo primer parámetro es \lstinline|[i32]|. En la segunda llamada ocurre lo contrario, ya que el literal de slice \lstinline|[0, 1, 2]| tiene tipo \lstinline|[i32]|, que solo unifica con la segunda firma. Ambas llamadas comparten el mismo nombre de función en el código fuente, pero el compilador selecciona una implementación distinta en cada punto de llamada sin ninguna anotación adicional por parte del programador.

\subsection{Detección de un error de tipos en tiempo de compilación}

El sistema de tipos rechaza en tiempo de compilación cualquier programa cuyas restricciones no puedan unificarse, incluso si la operación tendría sentido en un lenguaje de tipado dinámico. Considérese la siguiente declaración incorrecta.

\begin{lstlisting}[language={}, caption=Declaración que produce un error de unificación entre qbool e i32]
combinar :: qbool -> i32
combinar x = x + 1
\end{lstlisting}

La aplicación de \lstinline|+| exige, según su firma \lstinline|i32 -> i32 -> i32|, que su primer operando tenga tipo \lstinline|i32|, de modo que la generación de restricciones añade $\tau_x = \mathtt{i32}$. Sin embargo, la anotación de la declaración ya fijó $\tau_x = \mathtt{qbool}$ mediante la regla de inferencia de declaraciones. Al intentar unificar \lstinline|qbool| con \lstinline|i32| (\autoref{sec:unificacion}), ambos son tipos \lstinline|Builtin| con identificadores distintos, por lo que la unificación falla y el compilador emite un error señalando el span de \lstinline|x| dentro de la expresión \lstinline|x + 1|, junto con los dos tipos incompatibles. Este tipo de error, trivial de introducir por accidente en una biblioteca de tipado dinámico como Qiskit o Cirq (\autoref{sec:estadodelarte}), queda descartado en Qalc antes de que el programa llegue a generarse como circuito.
